\section{Paired paths and the cost of leaving a block language}
\label{sec:paired-paths}

This section gives paper proofs, not new Lean declarations. Paired parity
words and their concatenation are classical tools
\cite{garner1985,eliatucker2015}. Our purpose is to give a small, explicit
instance with positive integer states, unit repairs, a different final
parent, and quantified exit costs. Neither the mechanism nor the first
construction of a smaller source before descent is claimed as new;
the latter already has substantially stronger constructions in
\cite[Theorem~3]{weaver21}. The open issue here is coverage of an already
specified integer orbit.

\subsection{An exact block relation}

For a positive-exponent word $w=(a_1,\ldots,a_k)$, write
$A(w)=\sum_i a_i$ and
\[
 F_w(x)=\frac{3^k x+C_w}{2^{A(w)}},\qquad
 C_{()}=0,\qquad C_{w\mathbin{\|}(a)}=3C_w+2^{A(w)}.
\]
The function $F_w$ is an affine expression; it equals $S^k$ only on the
exact valuation cylinder of $w$. We use the exact-cylinder criterion
established in v6 \cite{borgatta_v6}: for a positive odd integer $x$,
the word $w$ is realized precisely when $F_w(x)$ is an odd integer.
For completeness, the reverse implication follows by induction from
$C_w=3^{k-1}+2^{a_1}C_{(a_2,\ldots,a_k)}$: oddness of the final quotient
forces the first valuation to be $a_1$, after which the same argument
applies to the remaining word. All actual intermediate states are
positive because the initial state is positive.

\begin{lemma}[Translation of exact blocks]\label{pp:translation}
Let $W,V$ have the same length $k\ge1$ and total exponent $A$.
For an even integer $c$,
\[
 F_V(x+c)=F_W(x)+c
 \quad\Longleftrightarrow\quad
 C_V-C_W=c(2^A-3^k).
\]
If this identity holds, $x$ realizes $W$, and both $x+c$ and
$F_W(x)+c$ are positive, then $x+c$ realizes $V$ exactly.
\end{lemma}
\begin{proof}
Compare the affine constants. The second assertion follows because
$F_W(x)+c$ is then a positive odd integer and the exact-cylinder
criterion applies to the positive odd input $x+c$.
\end{proof}

The translation $H(x)=x-2$ gives two rules:
\begin{equation}\label{pp:rules}
 W_1=(1,2)\longmapsto V_1=(2,1),\qquad
 W_2=(2,1,2)\longmapsto V_2=(1,1,3).
\end{equation}
Their affine constants and cylinders are
\[
\begin{array}{c|c|c|c|c}
 W&V&(3^k,2^A)&(C_W,C_V)&\text{source cylinders }W,V\\ \hline
 (1,2)&(2,1)&(9,8)&(5,7)&11,9\pmod {16}\\
 (2,1,2)&(1,1,3)&(27,32)&(29,19)&57,55\pmod {64}.
\end{array}
\]
In each row $F_V(x-2)=F_W(x)-2$. The least positive inputs
in these cylinders exceed $2$, and the block endpoints also exceed $2$.
Consequently any actual concatenation $p\in\{W_1,W_2\}^*$ starting
at $b>2$ is followed by an actual concatenation $\sigma(p)$ starting
at $b-2$, with a difference of $2$ at every block boundary.
The substitution $\sigma$ is defined by \eqref{pp:rules}.
This is intertwining between different affine maps. It does not require
one companion map to commute with every observed block. Indeed, the
fixed points of $F_{W_1}$ and $F_{W_2}$ are $-5$ and $29/5$.
The second block has gain $27/32<1$; the entire block language is not
an expansive language.

\subsection{An attributed terminal word and a mixed family}

The binary word used below comes from Cassian Mahler's manuscript
\cite[Section~4, Equation~(4.4)]{mahler152}. Its broader search and
survival certificates are not inputs to the following direct proof.
Set
\begin{align*}
 w_*={}&(1,1,2,1,2,1,2,1,2,1,2,1,2,2,2,2,1,2,1,2,2,1),\\
 g_*={}&(4,\underbrace{1,\ldots,1}_{15},2,2,1,4,3,3).
\end{align*}
Both have length $22$ and total exponent $34$. The recursive definition
of the affine constant gives
\[
 C_{w_*}=118309227953,\qquad C_{g_*}=181071347171,
 \qquad C_{g_*}-C_{w_*}=2\cdot3^{22}.
\]
Thus $F_{g_*}(x-2)=F_{w_*}(x)$. These integers provide a direct finite
check of the displayed identity; they do not certify the other results
of \cite{mahler152}.

\begin{proposition}[A mixed paired family]\label{pp:mixed-family}
Suppose a positive odd integer $b$ realizes $p w_*$, where
$p\in\{W_1,W_2\}^*$. Then $b-2$ realizes $\sigma(p)g_*$ and the
two paths have the same endpoint, length, total exponent, and gain.
Their last parents are different. For a unit input $b$, a unit source
for the alternative path is given by the following table; $\Gamma$
denotes its gain divided by the observed gain:
\[
\begin{array}{c|c|c}
\text{condition}&\text{alternative unit source}&\Gamma\\ \hline
b\equiv1\pmod3&b-2&1\\
b\equiv2\pmod3,\ p\text{ begins with }W_1&(3b-5)/4&4/3\\
b\equiv2\pmod3,\ p\text{ begins with }W_2&(3b-5)/2&2/3\\
b\equiv2\pmod3,\ p=()&(3b-5)/16&16/3.
\end{array}
\]
The last three paths have one fewer odd step.
\end{proposition}
\begin{proof}
The seed cannot be $1$: its valuation word consists of $2$'s, whereas
$p w_*$ contains a $1$. Hence $b>2$, so the companion seed is positive.
Iterate Lemma~\ref{pp:translation} along $p$, then use the displayed
terminal identity. Exact oddness of its endpoint proves the legality
of $g_*$. The observed last exponent is $1$, whereas the alternative
last exponent is $3$. Equal last parents would have the same valuation,
so the parents are different.

If $b\equiv1\pmod3$, the source $b-2$ is a unit. If
$b\equiv2\pmod3$, it is a positive odd multiple of $3$.
Delete its first actual step. That step has exponent $f=2,1,4$
in the three respective cases, and its endpoint is
$(3b-5)/2^f$. This endpoint is a positive odd unit. Deleting the
step multiplies the gain by $2^f/3$ and leaves the last parent
unchanged. The remaining length is at least $21$.
\end{proof}

The third unit source in the table can exceed $b$. It is not a
unit descent certificate. A useful subclass avoids this issue. Put
\[
 A=W_1^4,\qquad B=W_1^3 W_2.
\]
These blocks form a prefix code, have lengths $8,9$ and exponent totals
$12,14$, and are not powers of one common word. Every nonempty prefix
of either block has gain greater than $1$. For the last three letters
of $B$, the smallest intermediate gain is
$(9/8)^3(3/4)=2187/2048>1$; its total gain is $19683/16384>1$.
The exponent totals of the prefixes of $w_*$ are
\[
 (1,2,4,5,7,8,10,11,13,14,16,17,19,21,23,25,26,28,29,31,33,34).
\]
Each is below the corresponding $i\log_2 3$; their smallest gain is
$3^{21}/2^{33}>1$. Hence every actual word in
$\{A,B\}^*w_*$ stays above its initial source at every nonempty
prefix, while Proposition~\ref{pp:mixed-family} supplies a smaller
unit source with the same future. This includes the empty macro prefix.

These are identities on proper cylinders, not coverage of a supplied
infinite orbit. The relative odd Haar measure of all cylinders
$\{W_1,W_2\}^*w_*$ is only
\[
 \sum_{r\ge0}(2^{-3}+2^{-5})^r2^{-34}
 =\frac{2^{-29}}{27}.
\]
They are disjoint because $w_*$ starts with $11$, whereas the two
blocks start with $12$ or $212$. Nor does a merger alone certify a
first arrival. The first-arrival and prefix-free guards must be checked
when the earlier mass budgets are applied; for a nonperiodic target,
an actual arrival is automatically its unique arrival from that source.

\subsection{A minimum divergent seed must leave the language}

\begin{proposition}[A necessary exit; no analytic input]
\label{pp:minimum-exit}
If a divergent positive odd unit orbit exists, let $b$ be its least
unit starting value. Its infinite valuation word does not belong to
$\{W_1,W_2\}^{\mathbb N}$.
\end{proposition}
\begin{proof}
The seed satisfies $b>1$. Its first valuation must be $1$, since a
valuation at least $2$ gives the smaller divergent unit
$S(b)\le(3b+1)/4<b$. If its word were an infinite block concatenation,
its first block would therefore be $W_1$.

At block boundaries the actual orbit of $b-2$ differs from that of $b$
by $2$. Those values tend to infinity. An orbit of a deterministic map
on the positive integers that has an unbounded subsequence cannot
repeat a state, and therefore also tends to infinity. Thus $b-2$
diverges. If it is a unit, this contradicts minimality. Otherwise its
first valuation is the first alternative valuation of $V_1$, namely
$2$, and $S(b-2)=(3b-5)/4$ is a divergent unit smaller than $b$.
This is again a contradiction.
\end{proof}

The proposition requires neither a merger nor an atom-floor theorem.
It constrains the \emph{minimum} divergent unit seed, not every
divergent seed. It leaves open all exits not covered by a further
certificate. Elia and Tucker's counterexamples to a specific proposed
classification of paired stems \cite{eliatucker2015} are a relevant
warning: a collection of valid local rules does not establish their
completeness.

\subsection{Local costs and subsequent forced contraction}

At a block boundary write the pair as $(x,x-2)$, and let $x'$ be the
next observed state. If the observed exponent is $e\ge2$, the
companion exponent is exactly $1$ and the new relation is
$y'=2^{e-1}x'-3$. If $e=1$, the companion exponent is
$f=1+\nu_2(x'-3)\ge2$ and $y'=(x'-3)/2^{f-1}$.
These follow directly by subtracting $6$ from $3x+1$.
After the prefix $21$ the relation is $y=2x-5$.
An exit $13$ or $213$ consequently produces
\[
 y'=\operatorname{oddpart}(z-1),\qquad z=x' >1.
\]
Here $z=1$ is incompatible with an integral preceding state of
valuation $3$ in either case.

\begin{lemma}[A forced run]\label{pp:forced-two}
If $z=2^t u+1$, where $u$ is positive odd and $t\ge1$, then the next
$r=\lfloor(t-1)/2\rfloor$ valuations are all $2$ and
\[
 S^r(z)=(3/4)^r(z-1)+1.
\]
\end{lemma}
\begin{proof}
For $s\ge3$,
$3(2^s u+1)+1=4(3\cdot2^{s-2}u+1)$, with the factor in parentheses
odd. Iterate while the remaining exponent of $2$ is at least $3$.
\end{proof}

\begin{proposition}[Local cost or a drop below the reference seed]
\label{pp:buffer-cost}
Let $b>1$ be a reference value and suppose an observed path reaches
$z=2^t u+1$. Either one of $z,S(z),\ldots,S^r(z)$ lies below $b$,
where $r$ is as in Lemma~\ref{pp:forced-two}, or
\begin{equation}\label{pp:t-charge}
 t\le 2+\frac{2\log((z-1)/(b-1))}{\log(4/3)}.
\end{equation}
In particular the inequality holds along the orbit of a minimum
divergent unit seed $b$ at an exit $13$ or $213$.
For an exit beginning with $11$, if the first companion valuation is
$f\ge7$, the corresponding bound is
\begin{equation}\label{pp:f-charge}
 f\le 8+\frac{2\log((z-1)/(b-1))}{\log(4/3)}
\end{equation}
unless the observed path subsequently falls below $b$.
\end{proposition}
\begin{proof}
If there is no drop, the last forced state gives
$(3/4)^r(z-1)+1\ge b$. Hence
$r\log(4/3)\le\log((z-1)/(b-1))$, and $t\le2r+2$ proves
\eqref{pp:t-charge}. All later states of a minimum divergent unit
orbit are units at least $b$.

For the last assertion put $u=S(x-2)$, where $x$ is the boundary
state. From $3(x-2)+1=2^f u$, the first observed states are
\[
 x_1=3+2^{f-1}u,\quad x_2=5+3\cdot2^{f-2}u,\quad
 z=x_3=1+9\cdot2^{f-6}u.
\]
For $f\ge7$ the three observed exponents are exactly $1,1,4$.
Apply \eqref{pp:t-charge} with $t=f-6$.
\end{proof}

The cost in these inequalities is controlled by the actual height
buffer $\log((z-1)/(b-1))$. It is not an $O(\log j)$ cost at time $j$.
The remaining cases $f=3,4,5,6$ are four parametric classes, with
unbounded integer states, not four finite numerical states.

\begin{proposition}[A local loss of the same companion]
\label{pp:collapse-family}
For every fixed $K\ge5$, there are arbitrarily large positive odd
units $b$ whose observed word begins with $W_1^K(1,3)$, all its
nonempty prefixes have gain greater than $1$, and the companion
$b-2$ reaches $1$ while the observed path reaches an arbitrarily
large value.
\end{proposition}
\begin{proof}
Choose arbitrarily large $t=6m$ satisfying
$3^{2K+3}\mid 2^{t+1}+7$. Such $m$ exist because $64$ has order
$3^{N-2}$ modulo $3^N$ on $1+9\mathbb Z/3^N\mathbb Z$, for $N\ge2$,
as follows from $\nu_3(64^h-1)=2+\nu_3(h)$. The target $-7/2$
belongs to that subgroup. Put
\[
 x=\frac{2^{t+4}+11}{9},\qquad
 b=\frac{8^{K+1}(2^{t+1}+7)}{9^{K+1}}-5.
\]
For sufficiently large $t$ these are positive odd integers, with
$b\equiv1\pmod3$, so both initial seeds are units. The identity
$x+5=(9/8)^K(b+5)$ and the exact block cylinders give
\[
 b\xrightarrow{W_1^K}x\xrightarrow{(1,3)}2^t+1,
 \qquad
 b-2\xrightarrow{V_1^K}x-2\xrightarrow{(2,t+2)}1.
\]
The last intermediate states are $(2^{t+3}+7)/3$ and
$(2^{t+2}-1)/3$, which verify the last valuations exactly.
The smallest gain on the final two observed letters is
$(9/8)^K(9/16)>1$ for $K\ge5$, since $9^6>16\cdot8^5$.
All preceding prefixes are also expansive.
\end{proof}

This is not a divergent orbit construction or a counterexample to
Collatz. For fixed $K$, the ratio $(2^t)/(b-1)$ stays bounded as
$t\to\infty$, while the following forced run has unbounded length.
It therefore eventually drops below $b$. The example rules out a
uniformly bounded local cost for keeping this particular companion;
it does not rule out paying the cost using later contraction.

The next required result is a composition theorem: at every exit,
either certify a drop below the original seed, or choose another
positive unit companion and control the accumulated costs. The local
bounds above do not prove that result, since later expansive blocks
can replenish the height buffer. No universal descent, no complete
exit classification, and no nonsummable allocation of alternative
paths follows from this section.
